\par\Needspace{6\baselineskip}\medskip
\hypertarget{native-127-article094}{}\pdfbookmark[2]{94.}{article-094}
\begin{center}94.\end{center}
Die Behandlung der Gleichungen 12 und $12^{\ast}$, wie sie in Art. 91, 92 und 93 erläutert wurde, beruht auf der Annahme, dass der Winkel $g$ nicht zu groß ist und jedenfalls unter der Grenze $66^0\,25'$ liegt, über die wir Tafel III nicht hinausführen. Ist diese Annahme nicht erfüllt, so erfordern die Gleichungen kein derart aufwendiges Verfahren. Sie lassen sich dann \textit{in unveränderter Form} stets sehr sicher und bequem durch Probieren lösen. \textit{Sicher}, weil man den Wert des Ausdrucks $\dfrac{2g-\sin 2g}{\sin g^3}\mathpunct{\raisebox{.45ex}{,}}$ in dem $2g$ selbstverständlich in Teilen des Radius auszudrücken ist, für größere Werte von $g$ mithilfe trigonometrischer Tafeln \textit{mit voller Genauigkeit} berechnen kann, was bei einem kleinen Winkel $g$ nicht möglich ist. \textit{Bequem}, weil man heliozentrische Orte mit so großem Abstand voneinander wohl kaum je zur Bestimmung einer noch völlig unbekannten Bahn verwenden wird und sich bei irgendeiner Kenntnis der Bahn ein Näherungswert für $g$ nahezu mühelos nach Gleichung 1 oder 3 in Art. 88 ergibt. Aus dem Näherungswert von $g$ erhält man schließlich stets mit wenigen Versuchen den korrigierten Wert, der Gleichung 12 oder $12^{\ast}$ mit jeder gewünschten Genauigkeit genügt. Wenn die beiden gegebenen heliozentrischen Orte mehr als einen vollen Umlauf umfassen, muss man übrigens bedenken, dass die exzentrische Anomalie ebenso viele volle Umläufe vollendet haben wird. Die Winkel $E'-E$ und $v'-v$ liegen daher \SourceLeaf{128}entweder beide zwischen 0 und $360^0$ oder beide zwischen entsprechenden Vielfachen eines Vollkreises, und folglich liegen $f$ und $g$ entweder beide zwischen 0 und $180^0$ oder zwischen entsprechenden Vielfachen eines Halbkreises. Wäre schließlich die Bahn völlig unbekannt, sodass es ungewiss bliebe, ob der Himmelskörper beim Übergang vom ersten Radiusvektor zum zweiten nur einen Teil eines Umlaufs zurückgelegt oder darüber hinaus einen oder mehrere volle Umläufe vollendet hat, so könnte unsere Aufgabe bisweilen mehrere verschiedene Lösungen zulassen. Wir gehen hier über diesen Fall hinweg, da er in der Praxis kaum je vorkommen wird.
